\section{模型 Scaling：架构与超参数选择}

从数据 Scaling 转向模型 Scaling，我们面临的是更实际的工程问题：哪种架构更好？用什么优化器？深度和宽度怎么配比？

\subsection{架构比较：Transformer vs. LSTM}

\begin{figure}[H]
\centering
\includegraphics[width=0.95\textwidth]{slides-images/slide-017.jpg}
\caption{Kaplan：Transformer vs. LSTM 的 Scaling Law 对比\protect\footnotemark}
\end{figure}
\footnotetext{Slides 第18页。}

Kaplan 的论文清楚地展示了 Transformer 与 LSTM 在 Scaling Law 上的差异。关键观察：

\begin{itemize}
    \item 两种架构的 Scaling Law \textbf{斜率相似}（即缩放行为类似）
    \item 但存在一个\textbf{常数因子的偏移}：LSTM 在计算效率上比 Transformer 差约 15 倍
    \item 这个差距在 log 尺度上是\textbf{恒定的}，不会随着规模扩大而缩小
\end{itemize}

\begin{knowledgebox}{Scaling Law 视角下的架构选择}
Scaling Law 提供了一种优雅的架构选择方法：训练多种架构的小模型，跨几个数量级的计算量，然后比较它们的 Scaling 曲线。\textbf{斜率}告诉你两种架构的``缩放效率''是否相同；\textbf{偏移量}告诉你在固定计算量下哪个更好。
\end{knowledgebox}

\subsection{更广泛的架构比较}

\begin{figure}[H]
\centering
\includegraphics[width=0.95\textwidth]{slides-images/slide-018.jpg}
\caption{多种架构与 Transformer 的 Scaling 对比（Elhage et al., Google）\protect\footnotemark}
\end{figure}
\footnotetext{Slides 第19页。}

Google 的研究者对大量架构进行了系统性的 Scaling 比较。红线代表各种替代架构，绿线代表 Transformer 基线。结果表明：

\begin{importantbox}{能可靠超越 Transformer 的改进}
在所有测试的架构变体中，只有两类改进能\textbf{强且可靠地}超越标准 Transformer：
\begin{enumerate}
    \item \textbf{Gated Linear Units (GLU)}：门控线性单元
    \item \textbf{Mixture of Experts (MoE)}：专家混合模型
\end{enumerate}
这恰好就是当今前沿模型（如 LLaMA、Gemini）中实际使用的技术。Scaling Laws 在几年前就给出了正确的方向指引。
\end{importantbox}

\subsection{优化器选择：Adam vs. SGD}

\begin{figure}[H]
\centering
\includegraphics[width=0.95\textwidth]{slides-images/slide-019.jpg}
\caption{Adam vs. SGD 的 Scaling Law 对比（来自 Hestness et al.）\protect\footnotemark}
\end{figure}
\footnotetext{Slides 第20页。}

类似的分析可以应用于优化器选择。Hestness 等人比较了 SGD 和 Adam 优化器，发现二者之间也存在\textbf{常数因子的计算效率差距}，Adam 在相同计算量下能达到更低的损失。

\subsection{深度 vs. 宽度：Aspect Ratio}

\begin{figure}[H]
\centering
\includegraphics[width=0.95\textwidth]{slides-images/slide-020.jpg}
\caption{Kaplan：不同层数的 Scaling Law 对比\protect\footnotemark}
\end{figure}
\footnotetext{Slides 第21页。}

Kaplan 论文中关于深度与宽度比例的分析很有意思：
\begin{itemize}
    \item 1 层的模型明显很差
    \item 但从 2 层开始，不同层数的 Scaling 曲线非常接近
    \item 存在一个``很宽的最优盆地''：宽度/深度比在 4--16 之间都接近最优
\end{itemize}

\begin{warningbox}{不是所有参数都相同}
在做参数缩放分析时，一个容易犯的错误是将 embedding 层参数也计入总参数量。Kaplan 发现：
\begin{itemize}
    \item 如果包含 embedding 参数，Scaling 曲线会在某处``弯曲''，不再保持清晰的 log-log 线性
    \item 如果只考虑非 embedding 参数，曲线干净得多
\end{itemize}
原因是 embedding 层参数与 Transformer 层参数的``学习效率''不同。类似的问题也出现在 Mixture of Experts 中——稀疏激活的参数需要转换为``等效密集参数''。
\end{warningbox}

\begin{figure}[H]
\centering
\includegraphics[width=0.95\textwidth]{slides-images/slide-021.jpg}
\caption{Embedding 参数 vs. 非 Embedding 参数：后者的 Scaling 更干净\protect\footnotemark}
\end{figure}
\footnotetext{Slides 第22页。}

\subsection{超参数的跨尺度迁移}

\begin{figure}[H]
\centering
\includegraphics[width=0.95\textwidth]{slides-images/slide-022.jpg}
\caption{Kaplan 论文中的超参数 Scaling 分析：Aspect Ratio、前馈层维度比、注意力头维度\protect\footnotemark}
\end{figure}
\footnotetext{Slides 第23页。}

一个关键发现是：许多 Scaling Law 曲线在不同超参数设定下\textbf{斜率相同、不交叉}，只有常数偏移。这意味着：

\begin{importantbox}{超参数选择可以在小规模完成}
如果 Scaling 曲线对于不同超参数值只有常数偏移（斜率相同、不交叉），那么在\textbf{任意一个计算尺度}上选择最优超参数，结论就能迁移到所有更大的尺度。这大大降低了超参数搜索的成本。
\end{importantbox}

Kaplan 论文对多种超参数进行了这样的分析：Aspect Ratio（宽度/深度比）在不同规模的模型上表现一致；Feed-forward 维度比和注意力头维度也表现出类似的跨尺度稳定性。

\subsection{本章小结}

模型 Scaling Laws 为工程决策提供了强大的工具：通过在小规模上系统比较不同架构、优化器和超参数，可以可靠地预测大规模下的最优选择。关键洞察是大多数设计选择在 Scaling Law 中表现为常数偏移，使得小规模实验的结论可以直接迁移。

%% ========== Part 5: 批量大小与学习率 ==========

